% --- PREÁMBULO LIMPIO Y UNIFICADO DE LA EDITORIAL UD (preambulo_ud.tex) ---

% 1. CLASE DEL DOCUMENTO Y OPCIONES GLOBALES
% ¡Esta DEBE ser la primera línea de código no comentada!
%\documentclass[12pt,twoside,spanish]{book}

% 2. TIPOGRAFÍA Y ESTILOS DE LA EDITORIAL UD (REQUIERE LUALATEX)
\usepackage{fontspec} % Necesario para LuaLaTeX y la tipografía UD
\input{esConfigDir/esConfigFonts} % Definición de fuentes (lhseries, medseries, etc.)

\usepackage{geometry}
\geometry{
	paperheight=24.0cm,
	paperwidth=17.0cm,
	textwidth=12.5cm,
	tmargin=2.5cm,
	bmargin=2.5cm,
	left=2.5cm,
	right=2.0cm
	%   bindingoffset=1.1cm % Usar si se desea ajuste al encuadernar
}

% 3. ARCHIVOS DE ESTILO DE LA EDITORIAL UD
% Contiene \portadillauno, \espart, \nota, \fuente, y el estilo de página 'esstyle'.
%% Estilo LaTeX Espacios Editorial UD
%% Autores: Editores UDFJC
%% Versión:  v0.1 30MAR2024
%% Compilar con LuaLaTeX, i.e: lualatex main.tex
\usepackage[spanish]{babel}
\usepackage{csquotes}
\usepackage{caption}
\usepackage{titlesec}
\usepackage{titletoc}
\usepackage[titles]{tocloft}

%\setlength{\parindent}{1.25cm}
%\setlength\parindent{0pt}

\usepackage{xcolor}
\definecolor{MSBlack}{rgb}{0.0, 0.0, 0.0}
%\definecolor{MSBLack}{rgb}{.16,.16,.16}
\definecolor{MSLightBlack}{rgb}{.10,.10,.10}
\definecolor{MSCustomGray}{RGB}{144, 147, 150}
\definecolor{MSCustomBlackT}{RGB}{20, 20, 20}
\definecolor{MSCustomBlackST}{RGB}{58, 59, 66}

\newcommand\fontstit{\medseries}
\newcommand\fontschap{\regseries}
\newcommand\fontssec{\regseries}
\newcommand\fontssubsec{\regseries}
\newcommand\fontssubsubsec{\regseries}
\newcommand\fontsparagraph{\regseries\itshape}

% Small caps
\newcommand{\smallcaps}[1]{\textls[25]{\textsc{\MakeLowercase{#1}}}}

%configuracion para titulos de partes, capitulos, secciones, subsecciones, subsubsecciones
\newcommand{\chapterlabel}{\regseries Capítulo~\thechapter.}

% \titleformat{<command>}[<shape>]{<format>}{<label>}{<sep>}{<before-code>}[<after-code>]

\titleformat{\part}[display]{\color{MSBlack}\fontsize{25pt}{2.0mm}\selectfont\fontstit}{\thepart}{12pt}{\raggedright}
\titleformat{\chapter}[block]{\color{MSBlack}\fontsize{25pt}{2.5mm}\selectfont\fontschap}{\chapterlabel}{0.3em}{\raggedright}
\titleformat{\section}[block]{\color{MSBlack}\fontsize{16pt}{2.0mm}\selectfont\fontssec}{\thesection}{0.5em}{\raggedright}
\titleformat{\subsection}[block]{\color{MSBlack}\fontsize{13pt}{2.0mm}\selectfont\fontssubsec}{\thesubsection}{0.5em}{\raggedright}
\titleformat{\subsubsection}[block]{\color{black}\fontsize{12pt}{2.0mm}\selectfont\fontssubsubsec}{\thesubsubsection}{0.5em}{\raggedright}
\titleformat{\paragraph}[block]{\color{black}\fontsize{12pt}{2.0mm}\selectfont\fontsparagraph}{\theparagraph}{0.5em}{\raggedright}

%Espacios para titulos
\titlespacing{\part}{0cm}{10cm}{2.8cm}
\titlespacing{\chapter}{0cm}{1cm}{2.8cm}
\titlespacing{\section}{0cm}{7mm}{2mm}
\titlespacing{\subsection}{0cm}{5mm}{2mm} 
\titlespacing{\subsubsection}{0cm}{5mm}{0mm}
\titlespacing{\paragraph}{0cm}{5mm}{0mm}

% Configuración para apéndices
\makeatletter
\g@addto@macro{\appendix}{%
\titleformat{\chapter}[block]{\color{MSBlack}\fontsize{25pt}{2.5mm}\selectfont\fontschap}{Apéndice~\thechapter.}{0.3em}{\raggedright}
  \titlecontents{chapter}[0em]
    {\vspace{1.5mm}\medseries\normalsize}
    {Apéndice~\thecontentslabel.\enspace}
    {}
    {\hfill\small\lhseries\contentspage\vspace{2.5mm}}[]
}
\makeatother

%configuracion tabla de contenido
\titlecontents{part}[0em]
{\bf \sbseries\normalsize}%
{}% numbered sections formatting
{}% unnumbered sections formatting
{\titlerule*[1pc]{}\small\lhseries\contentspage \vspace{1.0mm}}[]%

\titlecontents{chapter}[0em]
{\vspace{1.5mm}\medseries\normalsize}%
{Capítulo~\thecontentslabel.\enspace}% numbered sections formatting
{}% unnumbered sections formatting
{\titlerule*[1pc]{}\footnotesize\lhseries\contentspage \vspace{1mm}}[]%

\titlecontents{section}[1.5em]
{\small\lhseries}%
{}% numbered sections formatting
{}% unnumbered sections formatting
{\titlerule*[1pc]{}\footnotesize\lhseries\contentspage \vspace{0.5mm}}[]%

\titlecontents{subsection}[3em]
{\lhseries}%
{}% numbered sections formatting
{}% unnumbered sections formatting
{\titlerule*[1pc]{}\small\lhseries\contentspage \vspace{0.5mm}}[]%

\usepackage{setspace}
\linespread{1.1}
\usepackage[skip=3.5pt, indent=20pt]{parskip}

\usepackage{fancyvrb}
\usepackage{fancyhdr}

\fancypagestyle{esstyle}{%
\fancyhead{}
  \fancyfoot{}
  \renewcommand{\headrulewidth}{0pt}
%   \newcommand{\runningtitle}{\fontsparagraph \scriptsize \texttitulocorto}
%  \fancyhf[HLE]{\lhseries\scriptsize Author's Name}
  \fancyhf[HLE]{\lhseries\scriptsize \loauthors}
%  \fancyhf[HRO]{\runningtitle}
  \fancyhf[HRO]{\fontsparagraph \scriptsize \texttitulocorto}
  \fancyfoot[LE]{{\fontsize{10.5}{1}\usefont{T1}{Roboto-TLF}{thin}{}\selectfont{E}}\hspace{-0.5mm}{\footnotesize\usefont{T1}{Inter-Sup}{thin}{}\selectfont{\reflectbox{S}}}{\usefont{T1}{Roboto-TLF}{thin}{}\selectfont{\thinspace|\thinspace}}{\lhseries\scriptsize\thepage}}

  \fancyfoot[RO]{{\lhseries\scriptsize\thepage}{\usefont{T1}{Roboto-TLF}{thin}{}\selectfont{\thinspace|\thinspace}}{\fontsize{10.5}{1}\usefont{T1}{Roboto-TLF}{thin}{}\selectfont{E}}\hspace{-0.5mm}{\footnotesize\usefont{T1}{Inter-Sup}{thin}{}\selectfont{\reflectbox{S}}}}
  \fancyfoot[C]{}
}

%https://www.latex-project.org/help/documentation/fntguide.pdf

\fancypagestyle{plain}{%
  \fancyhead{}
  \fancyfoot{}
  \renewcommand{\headrulewidth}{0pt}
  \fancyfoot[LE]{{\fontsize{10.5}{1}\usefont{T1}{Roboto-TLF}{thin}{}\selectfont{E}}\hspace{-0.5mm}{\footnotesize\usefont{T1}{Inter-Sup}{thin}{}\selectfont{\reflectbox{S}}}{\usefont{T1}{Roboto-TLF}{thin}{}\selectfont{\thinspace|\thinspace}}{\lhseries\scriptsize\thepage}}

  \fancyfoot[RO]{{\lhseries\scriptsize\thepage}{\usefont{T1}{Roboto-TLF}{thin}{}\selectfont{\thinspace|\thinspace}}{\fontsize{10.5}{1}\usefont{T1}{Roboto-TLF}{thin}{}\selectfont{E}}\hspace{-0.5mm}{\footnotesize\usefont{T1}{Inter-Sup}{thin}{}\selectfont{\reflectbox{S}}}}
  \fancyfoot[C]{}
}

\assignpagestyle{\chapter}{plain}

%\setcounter{table}{0}
\counterwithout{table}{chapter}
\counterwithout{figure}{chapter}
\counterwithout{equation}{chapter}
\captionsetup[figure]{labelfont=bf, labelsep=period, font=footnotesize, justification=centering, width = 0.8\textwidth}
\captionsetup[table]{labelfont=bf, labelsep=period, font=footnotesize, justification=centering, width = 0.8\textwidth}

%footnote format
\makeatletter
\renewcommand{\@makefntext}[1]{%
  \begin{list}{}{%
    \setlength{\labelwidth}{1.5em}%
    \setlength{\leftmargin}{\labelwidth}%
    \setlength{\labelsep}{3pt}%
    \setlength{\itemsep}{0pt}%
    \setlength{\partopsep}{0pt}%
    \setlength{\topsep}{0pt}%
    \footnotesize}%
  \item[\@makefnmark\hfil]#1%
  \end{list}%
}
\makeatother

\newcommand{\nota}[1]{\parbox{10cm}{\centering\scriptsize\color{MSBlack} {\bf Nota: }{#1}}}
\newcommand{\fuente}[1]{\parbox{10cm}{\centering\scriptsize\color{MSBlack} {\bf Fuente: }{#1}}}

\usepackage{wrapfig}

\newcommand{\espart}[6]{\part[Parte #1]{\protect\thispagestyle{empty}
%{
%\hspace{-2.75cm}\centering\includegraphics[width=17cm, height=8cm]{parte1_img}
%\begin{figure}[!h] %this figure will be at the right
\vspace{-14.0cm}
   \begin{figure}[ht!]\hspace{-2.9cm}\includegraphics[width=1.045\paperwidth, height=0.5\paperheight]{#2}\end{figure}
%\end{figure}
\flushleft
  \medseries{\Large\noindent Parte {#3}.}
  ~\\
  ~\\
  ~\\
  ~\\
  \lhseries{#4}
  ~\\
  ~\\
  ~\\
  ~\\
  \lhseries{\Large #5}
  ~\\
  ~\\
  ~\\
  ~\\
  \lhseries{\normalsize #6}
}}

%%%%%%%%%%%%%%%%%%


%%%%%%%%%% PORTADILLA 1 %%%%%%%%%%%%%%%
\newcommand{\portadillauno}[2]{
\vspace*{6.5cm}
\begin{flushleft}
  \linespread{3}
%  \color{MSCustomBlackT}\lsstyle
  \color{MSBlack}\lsstyle
  \fontsize{22pt}{3.0mm} \lhseries

  \leftskip=2mm
  #1   
  \vspace{2mm}
  \leftskip=-5mm

  %{\color{white}\colorbox{MSCustomGray}
    %{\parbox[][1.5cm][c]{15.2cm}{
        %\parbox{12.5cm}{\leftskip=7mm \fontsize{17pt}{2.5mm} \lhseries \vspace{1.0mm}
%          Consed ut escimos num sequid ut aut mos {\break} excerro et et, tempos quam, optat fuga.\vspace{0.5mm}
          %#2 \vspace{0.5mm}
        %}
      %}
    %}
  %}

\end{flushleft}}
%%%%%%%%%% FIN PORTADILLA 1 %%%%%%%%%%%%%%%

%%%%%%%%%% PORTADILLA 2 %%%%%%%%%%%%%%%
\newcommand{\portadillados}[2]{
\vspace*{6.5cm}
\begin{flushleft}
  \linespread{3}
%  \color{MSCustomBlackT}\lsstyle
  \color{MSBlack}\lsstyle
  \fontsize{22pt}{3.0mm} \lhseries

  \leftskip=2mm
  #1
   
  \vspace{2mm}
  \leftskip=-5mm

  {\parbox[][1.5cm][c]{15.2cm}{
      \parbox{12.5cm}{\leftskip=7mm \fontsize{17pt}{2.5mm} \lhseries \vspace{1.0mm}
%          Consed ut escimos num sequid ut aut mos {\break} excerro et et, tempos quam, optat fuga.\vspace{0.5mm}
          #2 \vspace{0.5mm}
      }
    }
  }
\end{flushleft}}
%%%%%%%%%% PORTADILLA 2 %%%%%%%%%%%%%%%

\makeatletter
\newcommand*\reftoanex[1]{\def\@currentlabelname{#1}}
\makeatother

%% COLOFON ESTILOS
\definecolor{MSGreyCOLOFON}{rgb}{.35,.35,.35}
\definecolor{MSGreyLATEX}{rgb}{.55,.55,.55}

%\usepackage{libertine}
\DeclareRobustCommand\myLaTeX{%
    L\kern-.3em%
    \raisebox{.4ex}{\scshape a}\kern-.13em%
    T\kern-.2em%
    \raisebox{-.5ex}{E}\kern-.13em%
    X%
}

\newcommand{\escolofon}[2]{
\thispagestyle{empty}
\vspace*{10.0cm}
\begin{tikzpicture}[remember picture,overlay]
  \node[text=gray,rotate=0,anchor=east, yshift=5mm] at (15.5,11.5){
    \includegraphics[width=5.0cm, height=4.0cm]{esIcons/logoESColofon_corner.png}
  };
\end{tikzpicture}

\vspace{-4.5cm}
\begin{center}
{\fontsize{7pt}{7pt}\selectfont
\color{MSGreyCOLOFON}
Este libro fue editado y diagramado en} {\fontsize{13pt}{13pt}\selectfont\color{MSGreyLATEX}{\myLaTeX}}
\vspace{3.5cm}

\includegraphics[scale=0.05]{esIcons/logoESColofon.png}
  
  \vspace{0.5cm}
  {\scriptsize
    \color{MSGreyCOLOFON}Este libro se término de imprimir en #1 de #2\\ en la Editorial UD. Bogotá, Colombia.}
\end{center}
}
%%%%%%%%%%%%%%%%%%

%\usepackage{showframe}% Similar al calderon, apagar al terminar de maquetar.
%\usepackage[axes,cam, width=190truemm, height=260truemm, pdftex,center]{crop}

% 4. CONFIGURACIÓN DE IDIOMA, CODIFICACIÓN Y TIPOGRAFÍA
% (Se mantienen tus paquetes, pero algunos se reordenan o eliminan si son redundantes)
\usepackage[spanish]{babel}
\usepackage{csquotes}
\usepackage{microtype} % Microtipografía avanzada
\usepackage{ragged2e} % Para control de alineación
%\usepackage[T1]{fontenc} % Se mantiene para compatibilidad, aunque fontspec es el principal.
%\usepackage{lmodern}
\usepackage{authblk} % Añadido para manejo de múltiples autores, si es necesario.
\decimalpoint % Para usar la coma como separador decimal.
\setlength\parindent{0pt} % Sin sangría al inicio de párrafos, como pide la UD.
\newcommand{\normalfonts}{\fontsize{12pt}{14.4pt}\selectfont} 
\let\normalsize\normalfonts
% 5. MATEMÁTICAS Y FÍSICA (TUS PAQUETES ORIGINALES)
\usepackage{amsmath}
\usepackage{amssymb}
\usepackage{amsfonts}
\usepackage{amsthm}
\usepackage{bm}
\usepackage{physics}
\usepackage{fontawesome}
\usepackage{fp}
\usepackage{calculus}
\usepackage{calculator}
\usepackage{fancyvrb}
\usepackage{mathtools}
\usepackage{bigints}
\usepackage{fontawesome}
% 6. GRÁFICOS E IMÁGENES (TUS PAQUETES ORIGINALES)
\usepackage{graphicx}
\usepackage{epstopdf}
\usepackage{auto-pst-pdf}
\usepackage[export]{adjustbox}
\graphicspath{{./Fig}{./Img}{./Imgn}}

% Paquetes de PSTricks
\usepackage{pstricks}
\usepackage{pstricks-add}
\usepackage{pst-node}
\usepackage{pst-3d}
\usepackage{pst-math}
\usepackage{pst-coil}
\usepackage{pst-3dplot}
\usepackage{pst-solides3d}

% 7. COLOR Y LAYOUT
\usepackage{xcolor}
\definecolor{orcidlogocol}{HTML}{A6CE39}
\usepackage{academicons}
\usepackage{tcolorbox}
\tcbuselibrary{breakable,skins}
\usepackage{multicol}
\usepackage{esvect}
\usepackage{relsize}
\usepackage{pifont}
\usepackage{etoolbox}
\newtcolorbox{graybox}{
	enhanced,
	boxrule=0pt, % Sin borde
	colback=gray!60!white, % Fondo gris oscuro (ajuste el porcentaje si es necesario)
	colupper=white, % Texto en blanco
	sharp corners, % Esquinas cuadradas
	boxsep=5pt, % Relleno interno
	left=1em, right=0.5em, top=0.5em, bottom=0.5em, % Más relleno
	nobeforeafter, % Asegura que no haya saltos de línea automáticos
}

% 8. FLOATS Y CAPTIONS
\usepackage{float}
\usepackage{placeins}
\usepackage{caption}
\usepackage{wrapfig}
\usepackage{array}
\usepackage{subcaption}
% Configuración de Leyendas (Captions) de Figuras y Tablas
\captionsetup[figure]{labelfont=bf, labelsep=period, font=footnotesize, justification=centering, width = 0.8\textwidth}
\captionsetup[table]{labelfont=bf, labelsep=period, font=footnotesize, justification=centering, width = 0.8\textwidth}
\usepackage{enumitem}
\usepackage{listings}
\usepackage{needspace}

% 9. BIBLIOGRAFÍA, HIPERVÍNCULOS E ÍNDICES
% NOTA: La plantilla UD usa APA, se sobreescribe tu estilo `authoryear` para que use el estilo de la editorial.
%\usepackage{hyperref}
\usepackage[hidelinks]{hyperref}
%\usepackage[
%backend=biber,
%style=apa, % Estilo de la Editorial UD (sobreescribe tu authoryear)
%]{biblatex}
\usepackage[
backend=biber,
style=ieee, % <-- CAMBIA DE APA A IEEE
sorting=none % Opcional: para que las referencias aparezcan en orden de citación
]{biblatex}
\addbibresource{Referencias/references.bib}
% *** INICIO DE LA CORRECCIÓN IEEE para idioma español (Solución para la "y") ***
% *** SOLUCIÓN IEEE: FORZAR CONJUNCIÓN EN INGLÉS ***
\DefineBibliographyStrings{spanish}{%
	and = {and}%
}
\DefineBibliographyStrings{english}{%
	andothers = {et\addabbrvspace al\adddot}
}

\DeclareDelimFormat{finalnamedelim}{%
	\ifnumgreater{\value{liststop}}{2}{\finalandcomma\addspace}{}%
	\bibstring{and}\addspace}
% *** FIN DE LA SOLUCIÓN ***
% *** FIN DE LA CORRECCIÓN IEEE ***
\usepackage{makeidx}
\usepackage{subfiles}

% REDEFINICION CITE/YEAR (Tu comando original, mantenido)(Cuando necesite APA solo descomentar)
%\DeclareCiteCommand{\citeyear}
%{\usebibmacro{prenote}}
%{\bibhyperref{\printfield{year}}}
%{\multicitedelim}
%{\usebibmacro{postnote}}

% Ajustes de estilo APA de la plantilla UD (Necesario para el formato editorial)
%\DefineBibliographyStrings{spanish}{andothers={\itshape{et~al}\adddot}}
%\makeatletter
%\renewcommand*{\apablx@ifrevnameappcomma}[2]{#2}
%\makeatother
%\DefineBibliographyExtras{spanish}{%	\let\finalandcomma=\empty}
%\DeclareDelimFormat[bib,textcite]{finalnamedelim}{%
	%\ifnumgreater{\value{liststop}}{2}{}{}%
	%\addspace\bibstring{and}\space}
%\DeclareDelimFormat[parencite]{finalnamedelim}{\addspace{y}\space}
%\setcounter{biburlbreakpenalty}{20}

% 10. COMANDOS PERSONALIZADOS Y ENTORNOS (TUS DEFINICIONES ORIGINALES)
\newcommand{\ds}{\displaystyle}
\newcommand{\bi}{\int}
\newcommand{\la}{\langle}
\newcommand{\ra}{\rangle}
\newcommand{\ei}[1]{\mathbf{e}_{#1}}
\newcommand{\ptc}{\tcbset{colback=gray!10!white,colframe=black!85!blue}\begin{tcolorbox}[breakable,enhanced,title=\bf Parametrizando el sólido]}
\newcommand{\ptr}{\tcbset{colback=gray!10!white,colframe=black!85!blue}\begin{tcolorbox}[breakable,enhanced,title=\bf Parametrizando la región de Integración]}
\newcommand{\tc}{\begin{tcolorbox}[breakable,enhanced,title=\bf Instrucción en {\tt Mathematica}]}
\newcommand{\etc}{\end{tcolorbox}}
\newcommand{\ve}[1]{\mathbf{#1}}
\DeclareMathOperator{\arcsinh}{arcsinh}
\DeclareMathOperator{\arccosh}{arccosh}
\DeclareMathOperator{\arctanh}{arctanh}
\DeclareMathOperator{\arccoth}{arccoth}
\DeclareMathOperator{\arcsech}{arcsech}
\DeclareMathOperator{\arccsch}{arccsch}
\DeclareMathOperator{\rot}{rot}
\newcommand{\cn}[3]{\pscircle(#1,#2){5pt}\rput(#1,#2){$#3$}}
\newcommand{\fig}{\begin{figure}}
\newcommand{\ef}[1]{\caption{#1}\end{figure}}
\newcommand{\vi}{\mathbf{i}}
\newcommand{\vj}{\mathbf{j}}
\newcommand{\vk}{\mathbf{k}}
\newcommand{\brake}{\\[4pt]} % Comando auxiliar para saltos
		
% 11. ENTORNOS DE TEOREMA Y REDEFINICIONES
\newtheorem{teorema}{Teorema}[section]
\newtheorem{prop}{Proposición}[section]
\newtheorem{defi}{Definición}[section]
\newtheorem{ejer}{Ejercicio}[section]
\newtheorem{ejemplo}{Ejemplo}[section]
\renewcommand{\proofname}{Demostración} % Redefinición para Español
\renewcommand*{\tablename}{\bf Tabla} % Estilo UD para nombres de tablas
\renewcommand*{\figurename}{\bf Figura} % Estilo UD para nombres de figuras
%\renewcommand*\contentsname{\lhseries Contenido} % Estilo UD para Tabla de Contenido
\counterwithout{footnote}{chapter}
\setcounter{tocdepth}{3}
\setcounter{secnumdepth}{3}
\renewcommand{\cfttoctitlefont}{\hfill\Huge}
\renewcommand{\cftdot}{}
		
% 12. SOLUCIÓN DE CONFLICTOS
\AtBeginDocument{\RenewCommandCopy\qty\qtySI} % Solución de conflicto physics
\renewcommand*\contentsname{\lhseries Contenido}
\makeindex
% 1. Restaurar el formato de numeración de secciones (ej: 1.1)
\renewcommand{\thesection}{\thechapter.\arabic{section}}

% 2. Restaurar el formato de numeración de subsecciones (ej: 1.1.1)
\renewcommand{\thesubsection}{\thesection.\arabic{subsection}}

% 3. Restaurar el formato de numeración de subsubsecciones (ej: 1.1.1.1)
\renewcommand{\thesubsubsection}{\thesubsection.\arabic{subsubsection}}

% (Bloque B: Restaura el FORMATO y añade el espacio)
\usepackage{titlesec} % <--- Asegúrate de que esté cargado antes de usar \titleformat

\titleformat{\section}[block]{\color{MSBlack}\fontsize{16pt}{2.0mm} \fontssec}
{\thesection.\quad}{0pt}{\raggedright}

\titleformat{\subsection}{\color{MSBlack}\fontsize{13pt}{2.0mm}\fontssubsec}
{\thesubsection.\quad}{12pt}{}

\titleformat{\subsubsection}[block]{\normalfont\bfseries\fontsize{13pt}{2.0mm}}{}{0pt}{\thesubsubsection\quad}

%\titleformat{\subsubsection}{\color{black}\bfseries\fontsize{13pt}{2.0mm}\fontssubsubsec}{\thesubsubsection.\quad}{12pt}{}

% --- En preambulo_ud.tex, en la Sección 12 ---
% SOLUCIÓN AL CONFLICTO DE COMANDO \fuente (Añadir si persiste el error)
\renewcommand{\fuente}[1]{\parbox{12.5cm}{\centering\footnotesize\color{MSBlack} {\bf Fuente: }{#1}}}
%REDEFINICIÓN DE NOTA: Usando tamaño absoluto (9pt)
\renewcommand{\nota}[1]{\parbox{12.5cm}{\centering\color{MSBlack} \fontsize{9pt}{11pt}\selectfont {\bf Nota: }{#1}}}

%***Corrección indice**
\makeatletter
% Restaura la numeración y el formato de los capítulos en el TOC
\renewcommand\l@chapter[2]{%
	\ifnum \c@tocdepth >\m@ne
	\vspace{0.5\baselineskip}%
	{\parindent \z@ \rightskip \@pnumwidth
		\leavevmode
		\bfseries
		\if@mainmatter \hskip 1.5em\fi % Indentación para capítulos
		\@tempdima \z@
		\advance\@tempdima by \the\c@chapter \relax % Fuerza el número del capítulo
		\hb@xt@\z@{\hfil#1}\hss\makebox[\@pnumwidth][r]{\makebox[1em][l]{#2}}\par}%
	\fi}

% Restaura la numeración y el formato de las secciones en el TOC
\renewcommand\l@chapter{\@dottedtocline{0}{0em}{1.5em}} % Nivel 0: Capítulo\renewcommand\l@section{\@dottedtocline{1}{1.5em}{2.3em}}

\renewcommand\l@section{\@dottedtocline{1}{1.5em}{2.3em}} % Nivel 1: Sección

% 3. Restaura la numeración y el formato de las subsecciones en el TOC
\renewcommand\l@subsection{\@dottedtocline{2}{3.8em}{3.2em}} % Nivel 2: Subsección

\renewcommand\l@subsubsection{\@dottedtocline{3}{5.8em}{4.1em}} % Nivel 3: Subsubsección

\makeatother

%***Cornizas personalizadas***

% 1. Activar el paquete para la limpieza de páginas en blanco
\usepackage{emptypage}

% 2. Cargar el paquete avanzado para la redefinición de cabeceras
\usepackage{fancyhdr}

% 3. Crear el nuevo estilo 'mis_cornisas'
\fancypagestyle{mis_cornisas}{
% Limpiar cabeceras y pies de página antes de definir
\fancyhf{} 
	
% Definir la línea negra inferior a la cabecera
\renewcommand{\headrulewidth}{0.4pt}
\renewcommand{\footrulewidth}{0pt} % Asegura que no haya línea en el pie
	
% Página Par (Izquierda): [Número de página] [Título del capítulo]
% \rightmark contiene el título del capítulo
\fancyhead[LE]{\thepage\hfill\nouppercase{\rightmark}} 
	
% Página Impar (Derecha): [Título del capítulo] [Número de página]
\fancyhead[RO]{\nouppercase{\leftmark}\hfill\thepage}
}

% 4. Definición para páginas iniciales de capítulo (estilo 'plain')
% Estas páginas no deben tener cornisa, solo el número centrado en el pie.
\fancypagestyle{plain}{
\fancyhf{} % Limpiar todo
\renewcommand{\headrulewidth}{0pt} % Quitar la línea
\fancyfoot[C]{\thepage} % Número de página centrado en el pie
}

% 5. Forzar la limpieza de página en blanco (para que no tengan nada)
\renewcommand{\cleardoublepage}{%
\clearpage\ifodd\value{page}\else\hbox{}\thispagestyle{empty}\newpage\fi%
}

% --- FIN DEL CÓDIGO DE CORNISA ---
		
% --- FIN DEL PREÁMBULO ---